% GENERATED FILE. Edit canonical lessons, then run npm run generate:books.
% canonical-source-hash: fnv1a64:341de6725f7cd9a8
% canonical-lessons: GU-C232-rog, GU-C232-cep, GU-C232-tandurast, GU-C232-kalji-levi, GU-C232-saju-thavu

\chapter{Staying healthy}
\label{ch:gu-staying-healthy}
\hlchaptermodality{232}

\begin{chapteropening}
\textbf{By the end of this chapter:} \emph{I can say how I stay healthy and ask about illness.}

Complete the last lesson of chapter 232: I can say how I stay healthy and ask about illness.
\end{chapteropening}

\section[\texorpdfstring{\gu{રોગ} (rog)}{રોગ (rog)}]{\gu{રોગ} (rog) — disease, illness}

\label{lesson:GU-C232-rog}



\begin{quote}
Before the new one: say the Gujarati for treatment, then the Gujarati for to be cured.
\end{quote}

\subsection*{You'll want to know: \gu{રોગ}}
\textbf{\gu{રોગ}} — \emph{rog} — \textquotedblleft{}disease, illness\textquotedblright{}.

\begin{quote}
\textbf{\gu{શું} \gu{આ} \gu{રોગ} \gu{ખરાબ} \gu{છે}?} — \emph{śũ ā rog kharāb che?} — \textquotedblleft{}Is this disease bad?\textquotedblright{}
\end{quote}

A masculine noun.

\begin{cousinweb}
\textbf{\gu{રોગ}} (\emph{rog}) is Sanskrit \emph{roga}, \textquotedblleft{}disease\textquotedblright{}, taken over whole as a learned loan, a tatsama (CDIAL 10829).
\end{cousinweb}

\begin{grammarlens}[title={What you've built}]
One more word for everyday life at A2.
\end{grammarlens}

\subsection*{Your turn}
\begin{itemize}
\raggedright
  \item \emph{Say it:} \emph{rog}
  \item \emph{Say it:} \emph{rog}, once more
  \item \emph{Say it:} \emph{śũ ā rog kharāb che?} — Is this disease bad?
  \item \emph{Recall:} say the Gujarati for treatment, then the Gujarati for to be cured, then say \emph{rog} again
\end{itemize}

\begin{tcolorbox}[breakable,colback=teal!4,colframe=teal!35!black,title={Before you move on}]
What does \gu{રોગ} mean? (\textquotedblleft{}Disease, illness\textquotedblright{}.) Say it once more.
\end{tcolorbox}
\section[\texorpdfstring{\gu{ચેપ} (cep)}{ચેપ (cep)}]{\gu{ચેપ} (cep) — infection, contagion}

\label{lesson:GU-C232-cep}



\begin{quote}
Before the new one: say the Gujarati for to be cured, then the Gujarati for disease.
\end{quote}

\subsection*{You'll want to know: \gu{ચેપ}}
\textbf{\gu{ચેપ}} — \emph{cep} — \textquotedblleft{}infection, contagion\textquotedblright{}.

\begin{quote}
\textbf{\gu{શું} \gu{તેને} \gu{ચેપ} \gu{છે}?} — \emph{śũ tene cep che?} — \textquotedblleft{}Does he have an infection?\textquotedblright{}
\end{quote}

A masculine noun.

\begin{cousinweb}
The origin of \textbf{\gu{ચેપ}} (\emph{cep}) is not settled: CDIAL lists no Gujarati form for it.
\end{cousinweb}

\begin{grammarlens}[title={What you've built}]
One more word for everyday life at A2.
\end{grammarlens}

\subsection*{Your turn}
\begin{itemize}
\raggedright
  \item \emph{Say it:} \emph{cep}
  \item \emph{Say it:} \emph{cep}, once more
  \item \emph{Say it:} \emph{śũ tene cep che?} — Does he have an infection?
  \item \emph{Recall:} say the Gujarati for to be cured, then the Gujarati for disease, then say \emph{cep} again
\end{itemize}

\begin{tcolorbox}[breakable,colback=teal!4,colframe=teal!35!black,title={Before you move on}]
What does \gu{ચેપ} mean? (\textquotedblleft{}Infection, contagion\textquotedblright{}.) Say it once more.
\end{tcolorbox}
\section[\texorpdfstring{\gu{તંદુરસ્ત} (tandurast)}{તંદુરસ્ત (tandurast)}]{\gu{તંદુરસ્ત} (tandurast) — healthy, fit}

\label{lesson:GU-C232-tandurast}



\begin{quote}
Before the new one: say the Gujarati for disease, then the Gujarati for infection.
\end{quote}

\subsection*{You'll want to know: \gu{તંદુરસ્ત}}
\textbf{\gu{તંદુરસ્ત}} — \emph{tandurast} — \textquotedblleft{}healthy, fit\textquotedblright{}.

\begin{quote}
\textbf{\gu{શું} \gu{તમારો} \gu{દીકરો} \gu{તંદુરસ્ત} \gu{છે}?} — \emph{śũ tamāro dīkro tandurast che?} — \textquotedblleft{}Is your son healthy?\textquotedblright{}
\end{quote}

It does not change for gender.

\begin{cousinweb}
\textbf{\gu{તંદુરસ્ત}} (\emph{tandurast}) is a loan, usually traced to Persian \emph{tandurust}, \textquotedblleft{}healthy\textquotedblright{}, built from \emph{tan}, \textquotedblleft{}body\textquotedblright{}, and \emph{durust}, \textquotedblleft{}sound\textquotedblright{}; EtymDB lists the Persian word.
\end{cousinweb}

\begin{grammarlens}[title={What you've built}]
One more word for everyday life at A2.
\end{grammarlens}

\subsection*{Your turn}
\begin{itemize}
\raggedright
  \item \emph{Say it:} \emph{tandurast}
  \item \emph{Say it:} \emph{tandurast}, once more
  \item \emph{Say it:} \emph{śũ tamāro dīkro tandurast che?} — Is your son healthy?
  \item \emph{Recall:} say the Gujarati for disease, then the Gujarati for infection, then say \emph{tandurast} again
\end{itemize}

\begin{tcolorbox}[breakable,colback=teal!4,colframe=teal!35!black,title={Before you move on}]
What does \gu{તંદુરસ્ત} mean? (\textquotedblleft{}Healthy, fit\textquotedblright{}.) Say it once more.
\end{tcolorbox}
\section[\texorpdfstring{\gu{કાળજી} \gu{લેવી} (kāḷjī levī)}{કાળજી લેવી (kāḷjī levī)}]{\gu{કાળજી} \gu{લેવી} (kāḷjī levī) — to take care, to look out for}

\label{lesson:GU-C232-kalji-levi}



\begin{quote}
Before the new one: say the Gujarati for infection, then the Gujarati for healthy.
\end{quote}

\subsection*{You'll want to know: \gu{કાળજી} \gu{લેવી}}
\textbf{\gu{કાળજી} \gu{લેવી}} — \emph{kāḷjī levī} — \textquotedblleft{}to take care, to look out for\textquotedblright{}.

\begin{quote}
\textbf{\gu{શું} \gu{તમે} \gu{તબિયતની} \gu{કાળજી} \gu{લો} \gu{છો}?} — \emph{śũ tame tabiyatnī kāḷjī lo cho?} — \textquotedblleft{}Do you take care of your health?\textquotedblright{}
\end{quote}

Word for word, to take care: \textbf{\gu{કાળજી}}, \emph{kāḷjī}, and \textbf{\gu{લેવું}}, \emph{levũ}.

\begin{cousinweb}
\textbf{\gu{કાળજી} \gu{લેવી}} (\emph{kāḷjī levī}) is built from \emph{kāḷjī}, \textquotedblleft{}care, anxiety\textquotedblright{}, and \emph{levũ}, \textquotedblleft{}to take\textquotedblright{}. CDIAL lists \emph{kāḷjī} under Sanskrit \emph{kāleyaka}, a part of the inner body near the liver, so care is felt in the heart (CDIAL 3103).
\end{cousinweb}

\begin{grammarlens}[title={What you've built}]
One more word for everyday life at A2.
\end{grammarlens}

\subsection*{Your turn}
\begin{itemize}
\raggedright
  \item \emph{Say it:} \emph{kāḷjī levī}
  \item \emph{Say it:} \emph{kāḷjī levī}, once more
  \item \emph{Say it:} \emph{śũ tame tabiyatnī kāḷjī lo cho?} — Do you take care of your health?
  \item \emph{Recall:} say the Gujarati for infection, then the Gujarati for healthy, then say \emph{kāḷjī levī} again
\end{itemize}

\begin{tcolorbox}[breakable,colback=teal!4,colframe=teal!35!black,title={Before you move on}]
What does \gu{કાળજી} \gu{લેવી} mean? (\textquotedblleft{}To take care, to look out for\textquotedblright{}.) Say it once more.
\end{tcolorbox}
\section[\texorpdfstring{\gu{સાજું} \gu{થવું} (sājũ thavũ)}{સાજું થવું (sājũ thavũ)}]{\gu{સાજું} \gu{થવું} (sājũ thavũ) — to get well, to recover}

\label{lesson:GU-C232-saju-thavu}



\begin{quote}
Before the new one: say the Gujarati for healthy, then the Gujarati for to take care.
\end{quote}

\subsection*{You'll want to know: \gu{સાજું} \gu{થવું}}
\textbf{\gu{સાજું} \gu{થવું}} — \emph{sājũ thavũ} — \textquotedblleft{}to get well, to recover\textquotedblright{}.

\begin{quote}
\textbf{\gu{શું} \gu{દર્દી} \gu{સાજો} \gu{થયો}?} — \emph{śũ dardī sājo thayo?} — \textquotedblleft{}Did the patient get well?\textquotedblright{}
\end{quote}

\textbf{\gu{સાજું}}, \emph{sājũ}, well, agrees with the person: \textbf{\gu{સાજો}}, \textbf{\gu{સાજી}}.

\begin{cousinweb}
\textbf{\gu{સાજું} \gu{થવું}} (\emph{sājũ thavũ}) is built from \emph{sājũ}, \textquotedblleft{}sound, healthy\textquotedblright{}, and \emph{thavũ}, \textquotedblleft{}to become\textquotedblright{}. CDIAL lists \emph{sājũ} under Sanskrit \emph{sajya}, \textquotedblleft{}strung (of a bow)\textquotedblright{}, so to be well is to be ready again (CDIAL 13095).
\end{cousinweb}

\begin{grammarlens}[title={What you've built}]
One more word for everyday life at A2.
\end{grammarlens}

\subsection*{Your turn}
\begin{itemize}
\raggedright
  \item \emph{Say it:} \emph{sājũ thavũ}
  \item \emph{Say it:} \emph{sājũ thavũ}, once more
  \item \emph{Say it:} \emph{śũ dardī sājo thayo?} — Did the patient get well?
  \item \emph{Recall:} say the Gujarati for healthy, then the Gujarati for to take care, then say \emph{sājũ thavũ} again
\end{itemize}

\begin{tcolorbox}[breakable,colback=teal!4,colframe=teal!35!black,title={Before you move on}]
What does \gu{સાજું} \gu{થવું} mean? (\textquotedblleft{}To get well, to recover\textquotedblright{}.) Say it once more.
\end{tcolorbox}
